\section{Results}
\label{sec:results}
\paragraph{Controlled reference agreement and incremental value.}
The final controls contain 77 no-load reference positives and 163
negatives. B4 and B5 have no observed disagreement with those labels. Their
H1 difference in correct acceptance per assigned case is 0.0
percentage points, with paired 97.5\% interval
[0.0, 0.0].
The empirical risk condition is satisfied, but the tied comparison does not
establish equivalence beyond these apparatus distributions. Positive geometry
alone (B3) also matches the nominal labels. No added-geometry accuracy benefit
is demonstrated. Relative to the no-load target, no-contact has 0 false
acceptances and 56 false rejections among 77 positives; its different
literal semantics remain intact (Table~\ref{tab:controls}).

\paragraph{Which distinctions survive targeted tests?}
In 8/8 registered representation pairs the complete
integration trace is bit-identical, while no-contact changes in 8.
B4 and B5 do not change. Accumulated no-load accepts
92 controls versus 77
for continuous time; the additional acceptances arise in the one-burst
intermittent group. The other nominal ablations do not establish an incremental
benefit on this controlled set. The registered branch subset gives
33 both-retain, 6 dependent and
9 unreachable states. Loaded-but-retaining redundant
supports separate realized load from dependence under this intervention.
No empty four-cell category is filled by adding favorable cases.

\paragraph{Natural tasks expose limits rather than a general accuracy gain.}
The historical no-contact exclusions yield 0/85 rescued scenes under no-load.
The new external results are source-specific (Table~\ref{tab:external}).
Fetch succeeds in 100/100 native episodes, while added retention has
0 definite passes, 52 failures and
48 indeterminate outcomes. Retained target damping torque
lies in the registered gray band; dropping torque would accept 48 cases but
would change the intended target. This is not evidence that the native expert
failed. The Shadow library yields 70 no-load passes,
29 failures and 1 indeterminate record;
1 record has invalid telemetry following a reproducible simulator
instability. Its native force-test successes are 66/100.
Withdrawal cannot test environmental dependence in this environment-free model.

\paragraph{Execution recovery and reference limits.}
5 of 720 Allegro candidate records fail the frozen
$10^{-10}$\,m geometry-consistency assertion. A documented fail-closed recovery
records them as invalid/indeterminate without widening the assertion, replacing
cases, changing physics or exposing future-use scores. Their native outcomes
are retained separately. This bookkeeping recovery was introduced after the
first execution failure and is not described as pre-test code. The 30-entry
probability reference samples per natural source remain separate from control
labels; unknown reference cases are not called verified failures or successes.

\paragraph{Candidate acceptance does not by itself prove utility.}
At K6, B5 accepts 13 and yields 10 usable
selections among all 120 scenes, with 107 abstentions.
B4 and B5 return the same candidate or abstention on 120/120 scenes;
their H4 difference is 0.0 points, paired 97.5\% interval
[0.0, 0.0].
No-contact and B5 differ in 0/120 K6 choices. Native-gate
selection provides 15 usable choices from
40 acceptances. These secondary contrasts concern the added
intervention-conditioned use, not improvements to the native program.
Unknown candidate futures make the finite-pool oracle a lower bound, so the
reported regret is an observed-pool diagnostic rather than true optimal regret.

\paragraph{Costs and scope of mechanical baselines.}
B4/B5 use the same nominal trajectories and no extra physical steps. Median
offline scoring calls on controls take 2.02 and 1.88\,ms,
excluding instrumentation, IO and simulation. These timings are
hardware-dependent and do not establish a runtime advantage. B6 and B7 are
reported on the fixed subset, with solver time and extra probe steps charged;
neither is relabeled as historical no-load ground truth. Appendix tables give
their applicability and outcomes. Native DGB execution and staging were
performed separately, but their step costs were not metered, so the recorded
audited-segment budget is not presented as an exhaustive benchmark cost.



\begin{table}[t]
\centering\small
\caption{Final no-load reference comparison on controlled apparatus. B0/B1/B2
have their own valid semantics; errors below are relative to the fixed $\NL$
reference, not accusations that their original definitions were implemented
incorrectly. U is indeterminate; invalid is separately reported in the text.}
\label{tab:controls}
\begin{tabular}{lrrrrr}\toprule
Method & Accepted & FP / 163 & FN / 77 & U & Accepted risk\\\midrule
B0 & 201 & 124 & 0 & 0 & 61.7\%\\
B1 & 182 & 105 & 0 & 0 & 57.7\%\\
B2 & 21 & 0 & 56 & 0 & 0.0\%\\
B3 & 77 & 0 & 0 & 0 & 0.0\%\\
B4\_normal & 77 & 0 & 0 & 0 & 0.0\%\\
B4 & 77 & 0 & 0 & 0 & 0.0\%\\
B5 & 77 & 0 & 0 & 0 & 0.0\%\\
A\_accumulated & 92 & 15 & 0 & 0 & 16.3\%\\
\bottomrule\end{tabular}

\end{table}

\begin{table}[t]
\centering\small
\caption{External source tasks remain separate. Allegro rows here use rank-1
candidates across 120 scenes, not 720 independent trials. Native external task
outcomes are not interchangeable with the added retention contract.}
\label{tab:external}
\begin{tabular}{lrrrrrr}\toprule
Source & N & Native & $\NC$ P & B4 P/F/U & B5 P/F/U & Invalid\\\midrule
Allegro rank-1 & 120 & 12 & 2 & 2/118/0 & 2/118/0 & 0\\
Fetch & 100 & 100 & 48 & 0/52/48 & 0/52/48 & 0\\
Shadow / DGB & 100 & 66 & 70 & 70/29/1 & 70/29/1 & 1\\
\bottomrule\end{tabular}

\end{table}

\begin{table}[t]
\centering\small
\caption{U2 on all 120 fresh Allegro scenes. Usable is independent future
retention of the selected candidate; abstentions remain in the denominator.
All methods pay the same complete nominal $K$-prefix simulation budget.}
\label{tab:utility}
\begin{tabular}{lrrrrr}\toprule
Method & K1 accept/use & K3 accept/use & K6 accept/use & K6 abstain & Regret\\\midrule
B0 & 12 / 2 & 34 / 11 & 40 / 15 & 80 & 0.283\\
B1 & 11 / 2 & 30 / 11 & 36 / 15 & 84 & 0.283\\
B2 & 2 / 2 & 8 / 6 & 13 / 10 & 107 & 0.325\\
B3 & 2 / 2 & 8 / 6 & 13 / 10 & 107 & 0.325\\
B4\_normal & 2 / 2 & 8 / 6 & 13 / 10 & 107 & 0.325\\
B4 & 2 / 2 & 8 / 6 & 13 / 10 & 107 & 0.325\\
B5 & 2 / 2 & 8 / 6 & 13 / 10 & 107 & 0.325\\
\bottomrule\end{tabular}

\end{table}

\begin{figure}[t]
\centering
\includegraphics[width=\linewidth]{risk_coverage_utility.pdf}
\caption{Registered control operating points and sealed candidate decisions.
Curve points change both normalized force and torque bounds together relative
to the fixed nominal reference. Overlapping B4/B5 curves must not be interpreted
as independent performance gains. The utility panel reports all-scene counts.}
\label{fig:curves}
\end{figure}
